% Standalone problem document, generated from the adjacent README.md.
% Compile with XeLaTeX. Mathematical content is maintained in Markdown.
\documentclass[11pt,a4paper]{article}
\usepackage[fontsize=10.5pt]{scrextend}
\usepackage[margin=22mm,headheight=15pt,headsep=9mm,footskip=11mm]{geometry}
\usepackage{fontspec}
\setmainfont{texgyrepagella}[Extension=.otf,UprightFont=*-regular,BoldFont=*-bold,ItalicFont=*-italic,BoldItalicFont=*-bolditalic]
\setsansfont{texgyreheros}[Extension=.otf,UprightFont=*-regular,BoldFont=*-bold,ItalicFont=*-italic,BoldItalicFont=*-bolditalic]
\setmonofont{texgyrecursor}[Extension=.otf,UprightFont=*-regular,BoldFont=*-bold,ItalicFont=*-italic,BoldItalicFont=*-bolditalic,Scale=MatchLowercase]
\usepackage{amsmath,amssymb}
\usepackage{unicode-math}
\setmathfont{texgyrepagella-math.otf}
\usepackage{xcolor,microtype,enumitem,needspace,fancyhdr}
\usepackage{bookmark}
\definecolor{ink}{HTML}{17344B}
\definecolor{accent}{HTML}{197D80}
\definecolor{muted}{HTML}{596773}
\hypersetup{colorlinks=true,linkcolor=accent,urlcolor=accent,pdftitle={KE-03 - Find
a near-largest nonnormal eigenvalue using few matrix-vector
products},pdfauthor={Open Problems in Numerical Linear Algebra}}
\urlstyle{same}
\setlength{\parindent}{0pt}
\setlength{\parskip}{5pt plus 1pt minus 1pt}
\linespread{1.04}
\setlength{\emergencystretch}{3em}
\widowpenalty=10000
\clubpenalty=10000
\displaywidowpenalty=10000
\allowdisplaybreaks[1]
\setlist{nosep,leftmargin=1.5em}
\providecommand{\tightlist}{\setlength{\itemsep}{0pt}\setlength{\parskip}{0pt}}
\providecommand{\pandocbounded}[1]{#1}
\pagestyle{fancy}
\fancyhf{}
\fancyhead[L]{\sffamily\footnotesize\color{muted}OPEN PROBLEMS IN NUMERICAL LINEAR ALGEBRA}
\fancyhead[R]{\sffamily\bfseries\footnotesize\color{accent}KE-03}
\fancyfoot[L]{\parbox[b]{0.9\textwidth}{\sffamily\fontsize{8}{10}\selectfont\color{muted}Editorial ratings; a bounded search cannot certify that a problem remains unsolved.\\Literature check: 2026-09-11}}
\fancyfoot[R]{\sffamily\footnotesize\color{muted}\thepage}
\renewcommand{\headrulewidth}{0.3pt}
\renewcommand{\headrule}{\hbox to\headwidth{\color{accent}\leaders\hrule height \headrulewidth\hfill}}
\makeatletter
\renewcommand{\section}{\Needspace{5\baselineskip}\@startsection{section}{1}{0pt}{12pt plus 2pt minus 2pt}{4pt}{\sffamily\bfseries\large\color{ink}}}
\renewcommand{\subsection}{\Needspace{4\baselineskip}\@startsection{subsection}{2}{0pt}{9pt plus 1pt minus 1pt}{3pt}{\sffamily\bfseries\normalsize\color{ink}}}
\makeatother
\setcounter{secnumdepth}{0}
\begin{document}
{\sffamily\small\color{muted}Eigenvalues and inverse problems\par}
\vspace{3pt}
{\raggedright\hyphenpenalty=10000\exhyphenpenalty=10000\sffamily\bfseries\fontsize{21}{24}\selectfont\color{ink}Find
a near-largest nonnormal eigenvalue using few matrix-vector
products\par}
\vspace{7pt}
{\sffamily\small\textbf{Difficulty:} challenging\quad\textbf{Importance:} interesting
to the community\par}
{\sffamily\small\bfseries\color{accent}Status: Solved\par}
\vspace{4pt}
{\color{accent}\hrule height 0.8pt}
\vspace{6pt}
\textbf{Topic:} nonsymmetric eigenvalue computation; query complexity

\textbf{Rating rationale:} Challenging reflects a query bound for
general nonnormal matrices, whose eigenvalues need not be controlled by
power iteration; community impact is the complexity of large
nonsymmetric eigenproblems.

\subsection{Resolution --- 2026-09-11}\label{resolution-2026-09-11}

\textbf{Affirmative resolution by Matthew J. Colbrook} (Department of
Applied Mathematics and Theoretical Physics, University of Cambridge).
See the
\href{https://github.com/ajt60gaibb/OpenProblemsInNLA/blob/main/eigenvalues-and-inverse-problems/KE-03/solution.md}{complete
manuscript}, \textbf{Theorem KE-03, sections 1--5}
(\href{https://github.com/ajt60gaibb/OpenProblemsInNLA/blob/main/eigenvalues-and-inverse-problems/KE-03/solution.pdf}{PDF}
·
\href{https://github.com/ajt60gaibb/OpenProblemsInNLA/blob/main/eigenvalues-and-inverse-problems/KE-03/solution.tex}{LaTeX}),
prepared 11 September 2026.

The algorithm uses \(O(\varepsilon^{-2}[1+\log(nK)])\) exact
matrix-vector queries, with success probability at least \(0.997\), for
every input in the displayed model. It supplies both eigenvalue-location
guarantees using the given condition bound \(K\) and finite exact
arithmetic between queries. The result bounds query count, not total
runtime, bit complexity or floating-point error.

The original proof draft was generated in a ChatGPT conversation. A
separate Codex agent independently verified the full proof and its match
to the exact target on 11 September 2026:
\href{https://github.com/ajt60gaibb/OpenProblemsInNLA/blob/main/references/colbrook-2026-09-11/verification/reviews/KE-03-review.md}{detailed
PASS review}. The review records a hash of the unchanged proof text.
This is independent agent verification, not external human peer review
or formal certification.
\href{https://github.com/ajt60gaibb/OpenProblemsInNLA/blob/main/references/colbrook-2026-09-11/README.md}{The
submission history and diagnostic record} preserve the initial solution
claim. The ratings above are historical, and the earlier literature
checks below are retained.

\subsection{Context and notation}\label{context-and-notation}

All norms are Euclidean vector norms or their induced matrix norms.

\subsection{Problem statement}\label{problem-statement}

An unknown diagonalizable \(A\in\mathbb C^{n\times n}\) is available
only through exact queries \(v\mapsto Av\). A number \(K\geq1\) is
supplied, with the promise that \(A=V\Lambda V^{-1}\) for some diagonal
\(\Lambda\) and \(\|V\|_2\|V^{-1}\|_2\leq K\). Assume
\(\rho(A)=\max_{\mu\in\sigma(A)}|\mu|>0\); neither the eigenvalues nor
\(V\) are supplied.

Do universal constants \(C,a,b>0\) and a randomized algorithm exist
that, for every such input and every \(\varepsilon\in(0,1/2)\), returns
a complex number \(z\) with probability at least \(0.99\) such that

\[
\text{there exists }\mu\in\sigma(A):
\quad |\mu|\geq(1-\varepsilon)\rho(A),
\qquad |z-\mu|\leq\varepsilon\rho(A),
\]

using at most \(C[1+\log(nK)]^a\varepsilon^{-b}\) matrix-vector queries?
The algorithm may perform finite exact arithmetic between queries; the
requested bound concerns query count. Products with \(A^*\) and
solutions of shifted linear systems are not additional oracle
operations. Supplying \(K\) and fixing a success probability make the
computational formulation explicit.

\newpage
\subsection{References}\label{references}

Amsel et al.,
\href{https://arxiv.org/html/2602.05394v3\#S3.SS5}{workshop report},
Problem 3.9. Shah, Srivastava, and Zeng,
\href{https://arxiv.org/pdf/2411.19926v1}{\emph{Sparse Pseudospectral
Shattering}, first version}, §1.3, Theorem 1.5 and its power-iteration
argument, supplies the related spectral-radius estimate cited by the
workshop.

\subsection{Earlier status check ---
2026-09-08}\label{earlier-status-check-2026-09-08}

Searches for \textit{nonnormal eigenvalue query complexity 2026} and
\textit{nonnormal spectral radius matrix-vector algorithm 2026}
found no resolution. The radius estimate controls a nonnegative scalar
and allows backward error; it does not locate a near-extremal eigenvalue
in the complex plane as required here. The
\href{https://arxiv.org/abs/2411.19926v3}{April 2026 third version}
replaces that application with a GMRES application (§6), so the
version-specific locator above is intentional. The workshop still poses
Problem 3.9 in its August update, separately from its resolved general
Ritz-compression question.

\subsection{Audit update --- 2026-09-10}\label{audit-update-2026-09-10}

Rechecked \href{https://arxiv.org/html/2602.05394v3}{workshop version
3}, Problem 3.9. Searches for nonnormal largest-eigenvalue matrix-vector
query bounds found no theorem or lower bound settling the displayed
question. Hermitian or singular-value oracle results do not answer it.
\end{document}
